\documentclass[10pt,a4paper]{article}
\usepackage[margin=2cm]{geometry}
\usepackage[T1]{fontenc}
\usepackage{lmodern}
\usepackage{microtype}
\usepackage{amsmath,amssymb}
\makeatletter\renewcommand\verbatim@font{\footnotesize\ttfamily}\makeatother
\sloppy
\usepackage{booktabs}
\usepackage{tabularx}
\usepackage{enumitem}
\usepackage[hidelinks]{hyperref}
\setlength{\emergencystretch}{3em}
\setlist{nosep}
\newlist{steps}{itemize}{1}
\setlist[steps]{label=$\square$, leftmargin=1.5em, itemsep=2pt}
\newcolumntype{Y}{>{\raggedright\arraybackslash}X}
\newcommand{\f}[1]{\texttt{#1}}

\title{BTP-I Implementation Plan\\[2pt]\large State UPI-adoption panel, descriptive results, report and slides}
\date{1 November 2026}
\begin{document}
\maketitle

\noindent\textbf{For agentic workers:} REQUIRED SUB-SKILL: use \f{superpowers:subagent-driven-development}
(recommended) or \f{superpowers:executing-plans} to implement this plan task by task. Each step is a
checkbox ($\square$); mark it done ($\blacksquare$) when complete.

\paragraph{Goal.} Build the quarterly state panel of UPI adoption (PhonePe Pulse, 2018Q1--2026Q2), validate it
against NPCI, produce the descriptive results, and set up the \LaTeX{} BTP-I report and Beamer slides.

\paragraph{Architecture.} Python scripts turn each raw source into one long interim table; one build script
merges them into \f{data/final/panel\_quarterly.csv} and stops on any failed check. R reads the final
panel and writes PDF figures and \f{.tex} tables, which the report and slides include.

\paragraph{Tech stack.} Python 3.12 (pandas, pytest) in \f{.venv}; R $\geq$ 4.3 with \f{renv}
(data.table, ggplot2, sf, modelsummary, testthat); pdflatex + BibTeX + natbib via \f{latexmk}; GNU Make; git.

\paragraph{Spec.} \f{docs/superpowers/specs/2026-11-01-upi-financial-inclusion-design.tex} (Sections 3--6).

\paragraph{Scope.} This plan covers only the BTP-I scope (spec Section~6). The FI data, FI indices and
Steps~1--3 (BTP-II) get a separate plan after the week-1 checks in Task~0. That plan cannot be exact until
we know which RBI indicators are quarterly.

\section*{Global constraints}
\begin{itemize}
\item PhonePe Pulse commit: \f{257bf38e3095020f9f516c4c72c10bdc348781b1} (branch \f{main}, 24 Sep 2026).
\item Period: \f{2018Q1} to \f{2026Q2}, 34 quarters. Period strings: quarterly \f{YYYYQn}, monthly \f{YYYY-MM}.
\item \f{state\_id}: lowercase ASCII snake case. Ladakh $\to$ \f{jammu\_kashmir}. Dadra \& Nagar Haveli and
      Daman \& Diu $\to$ \f{dnh\_dd}. 35 ids in total.
\item Main sample (30 ids): the 28 states + \f{delhi} + \f{jammu\_kashmir}. The other UTs are
      \f{andaman\_nicobar}, \f{chandigarh}, \f{dnh\_dd}, \f{lakshadweep}, \f{puducherry}.
\item Interim schema (every \f{data/interim/*.csv}): columns \f{state\_id,period,variable,value,source}, in this order.
\item Additive variables (counts, amounts, users, population) are summed over merged ids. Non-additive variables
      (CPI) take the value of the larger-population component.
\item Money: PhonePe \f{amount} is in rupees. Real values are in 2012 rupees: $real = nominal \times 100 / CPI$.
\item $\ln$ of a value $\le 0$ is \f{NaN}, never $-\infty$.
\item Outlier rule: a per-capita value above $10\times$ the median of all states in the same period fails the build.
      If real data trips this rule, stop and report it to the user. Do not change the threshold without approval.
\item Raw files are never edited. Each \f{data/raw/<source>/manifest.csv} has columns \f{file,url,retrieved,sha256,notes}.
\item Python modules run from the repo root as \f{PYTHONPATH=src/py .venv/bin/python -m btpdata.<module>}; each
      data module has a \f{main()} called under \f{if \_\_name\_\_ == "\_\_main\_\_"}.
\item All documents are \LaTeX{}. Citations use natbib + BibTeX (\f{plainnat}), not biblatex.
\item The student writes the report prose. Tasks~10--11 create the structure, outline comments, bibliography,
      figure and table includes. Agents do not write chapter prose unless the user asks.
\end{itemize}

\section*{Review focus}
\begin{enumerate}
\item State-name variants across sources (``Orissa'', ``Andaman And Nicobar Islands'', ``NCT of Delhi'', case and
      spaces) map to one id, and an unknown name stops the build and names the string. Tested in Task~1.
\item Merged ids: the per-capita value for \f{jammu\_kashmir} equals (J\&K + Ladakh sum) / (summed population), not
      a mean of two per-capita values. Tested in Tasks~2, 3, 5.
\item A zero or missing value (small UTs, early quarters) gives \f{NaN} logs and is counted in the coverage report;
      it does not stop the build. Tested in Task~5.
\item Population reference date: projections are as of 1 March, so the annual value is the Q1 value; 2026Q2
      needs the 2027 value. No population cell is empty in 2018Q1--2026Q2. Tested in Task~3.
\item A CPI base change inside the download (2012=100 to a newer base) is linked with the official factor, with
      no jump at the switch. Tested in Task~4.
\end{enumerate}

\section*{File map}
\begin{tabularx}{\textwidth}{@{}l Y@{}}
\toprule
File & Responsibility \\
\midrule
\f{pyproject.toml}, \f{requirements.txt} & pytest config (\f{pythonpath = ["src/py"]}), pinned deps \\
\f{src/py/btpdata/states.py}, \f{state\_crosswalk.csv} & name $\to$ \f{state\_id}; sample lists \\
\f{src/py/btpdata/periods.py} & period sequences \\
\f{src/py/btpdata/fetch\_phonepe.py}, \f{parse\_phonepe.py} & pinned sparse clone; JSON $\to$ \f{interim/phonepe.csv} \\
\f{src/py/btpdata/population.py} & \f{data/manual/population\_ncp2020.csv} $\to$ \f{interim/population.csv} \\
\f{src/py/btpdata/cpi.py} & MoSPI state CPI $\to$ \f{interim/cpi.csv} \\
\f{src/py/btpdata/npci.py}, \f{validation.py} & NPCI monthly $\to$ \f{interim/npci.csv}; \f{final/validation\_npci.csv} \\
\f{src/py/btpdata/boundaries.py} & state boundary GeoJSON $\to$ \f{interim/india\_states.geojson} with \f{state\_id} \\
\f{src/py/btpdata/build\_panels.py} & merge, derive, check $\to$ \f{final/panel\_quarterly.csv}, \f{coverage\_report.csv} \\
\f{src/R/lib/descriptives\_fns.R}, \f{src/R/01\_descriptives.R} & functions; figures and tables \\
\f{tests/py/}, \f{tests/R/} & tests and small fixtures \\
\f{Makefile} & \f{make test | data | analysis | report | all} \\
\f{thesis/btp1/}, \f{thesis/btp1-slides/} & report and slides \\
\bottomrule
\end{tabularx}

\section*{Parts (submission units)}
The work is built and submitted in three parts. Each part builds on the one before it, ends with a git tag,
and adds a progress report \f{thesis/progress/partN.tex} (what was built, the numbers, open issues).
\begin{description}
\item[Part 1, tag \f{part1}] Tasks 1--3: scaffold, state crosswalk, PhonePe panel, population.
      Report skeleton with Ch.~1--2 outlines (first half of Task~10).
\item[Part 2, tag \f{part2}] Tasks 4--7: CPI, quarterly panel build and checks, NPCI validation, boundaries.
      Needs the manual CPI and NPCI downloads.
\item[Part 3, tag \f{part3}] Tasks 8--11: R descriptives, Makefile end-to-end, full report structure, slides.
      Needs R installed.
\end{description}
Commits keep their real dates; the parts are not backdated.

% =====================================================================
\section{Task 0: Week-1 blocking checks (human)}
\begin{steps}
\item Email Prof.\ Pradhan: (a) PhonePe Pulse as the adoption proxy with NPCI validation (spec risk~1);
      (b) the two-tier design. Record his answer.
\item Get from the department BTP coordinator: the end-sem date, report format and page limit, template, and slide
      time (spec risk~9). Record them.
\item Install R and the system libraries for \f{sf}. The user runs:
      \f{! sudo apt install r-base r-base-dev libudunits2-dev libgdal-dev libgeos-dev libproj-dev}.
      Check: \f{Rscript -e 'cat(R.version.string)'} prints 4.3 or later.
\item Write the answers to \f{docs/decisions.tex} (a dated list) and commit:
      \f{git commit -m "docs: record week-1 decisions"}.
\end{steps}
If the professor rejects the proxy, stop and re-plan (spec risk~1 fallback).

% =====================================================================
\section{Task 1: Project scaffold and state crosswalk}
\textbf{Files.} Create \f{pyproject.toml}, \f{requirements.txt} (\f{pandas==2.2.3}, \f{pytest==8.3.3}),
\f{src/py/btpdata/\_\_init\_\_.py}, \f{states.py}, \f{state\_crosswalk.csv}, \f{tests/py/test\_states.py};
add \f{.venv/} to \f{.gitignore} (already present).

\textbf{Interfaces (produces).}
\f{normalise(raw: str) -> str} (lowercase, strip, ``\&''$\to$``and'', collapse spaces);
\f{to\_state\_id(raw: str) -> str}; \f{class UnknownStateError(ValueError)};
\f{MAIN\_SAMPLE: tuple[str, ...]} (30); \f{ALL\_IDS: tuple[str, ...]} (35).
The crosswalk CSV has columns \f{raw\_name,state\_id}, and \f{raw\_name} is already normalised.

\begin{steps}
\item Create the venv: \f{uv venv .venv \&\& uv pip install -r requirements.txt}.
\item Write the failing tests:
\begin{verbatim}
PHONEPE_36 = [...]  # the 36 names in PhonePe map/transaction/hover 2026/2.json, verbatim
@pytest.mark.parametrize("raw,sid", [
  ("Orissa","odisha"), ("ODISHA ","odisha"), ("Andaman And Nicobar Islands","andaman_nicobar"),
  ("andaman & nicobar islands","andaman_nicobar"), ("Jammu and Kashmir","jammu_kashmir"),
  ("ladakh","jammu_kashmir"), ("NCT of Delhi","delhi"), ("Pondicherry","puducherry"),
  ("dadra & nagar haveli & daman & diu","dnh_dd"), ("Dadra and Nagar Haveli","dnh_dd"),
  ("Daman and Diu","dnh_dd"), ("Uttaranchal","uttarakhand")])
def test_variants(raw, sid): assert to_state_id(raw) == sid
def test_unknown():
    with pytest.raises(UnknownStateError, match="Atlantis"): to_state_id("Atlantis")
def test_phonepe_names(): assert {to_state_id(n) for n in PHONEPE_36} == set(ALL_IDS)
def test_samples(): assert len(MAIN_SAMPLE)==30 and len(ALL_IDS)==35 and set(MAIN_SAMPLE)<=set(ALL_IDS)
\end{verbatim}
      The 36 PhonePe names (verified at the pinned commit): \f{andaman \& nicobar islands, andhra pradesh,
      arunachal pradesh, assam, bihar, chandigarh, chhattisgarh, dadra \& nagar haveli \& daman \& diu, delhi, goa,
      gujarat, haryana, himachal pradesh, jammu \& kashmir, jharkhand, karnataka, kerala, ladakh, lakshadweep,
      madhya pradesh, maharashtra, manipur, meghalaya, mizoram, nagaland, odisha, puducherry, punjab, rajasthan,
      sikkim, tamil nadu, telangana, tripura, uttar pradesh, uttarakhand, west bengal}.
\item Run \f{.venv/bin/pytest tests/py/test\_states.py -v}. Expected: FAIL (ImportError).
\item Implement \f{states.py} and fill \f{state\_crosswalk.csv} (one row for each variant above, plus each canonical name).
\item Run the tests. Expected: all PASS.
\item Commit: \f{git add pyproject.toml requirements.txt src/py tests/py \&\& git commit -m "feat: state crosswalk"}.
\end{steps}

% =====================================================================
\section{Task 2: PhonePe fetch and parse}
\textbf{Files.} Create \f{src/py/btpdata/periods.py}, \f{fetch\_phonepe.py}, \f{parse\_phonepe.py},
\f{tests/py/test\_phonepe.py}, and fixtures under \f{tests/py/fixtures/phonepe/} at the repo's relative
paths \f{data/map/<kind>/hover/country/india/2018/1.json}, with \f{<kind>} = \f{transaction}, \f{user}.

\textbf{Interfaces.} Consumes \f{to\_state\_id}. Produces:
\f{quarters(start: str = "2018Q1", end: str = "2026Q2") -> list[str]};
\f{PHONEPE\_COMMIT: str}; \f{fetch(dest: Path, commit: str = PHONEPE\_COMMIT) -> Path};
\f{parse\_transactions(path: Path, period: str) -> pd.DataFrame};
\f{parse\_users(path: Path, period: str) -> pd.DataFrame};
\f{build\_interim(repo: Path, out: Path, periods: list[str]) -> pd.DataFrame}.
Variables: \f{phonepe\_count}, \f{phonepe\_amount\_inr}, \f{phonepe\_users}; \f{source = "phonepe@<commit[:7]>"}.

\textbf{JSON layout (verified at the pinned commit).}
Transactions: \f{data.hoverDataList[] = \{name, metric:[\{type:"TOTAL", count, amount\}]\}}.
Users: \f{data.hoverData = \{<name>: \{registeredCount\}\}}. Path: \f{data/map/<kind>/hover/country/india/<year>/<q>.json}.

\begin{steps}
\item Write fixtures. Transactions 2018Q1: andhra pradesh (count 10862829, amount 14001071206.72507),
      jammu \& kashmir (100, 1000.0), ladakh (10, 50.0). Users: the same three names with 500, 40, 5.
\item Write the failing tests:
\begin{verbatim}
def test_quarters(): q = quarters(); assert len(q)==34 and q[0]=="2018Q1" and q[-1]=="2026Q2"
def test_parse_and_merge(tmp_path):
    df = build_interim(FIX, tmp_path/"p.csv", ["2018Q1"])
    assert list(df.columns) == ["state_id","period","variable","value","source"]
    v = df.set_index(["state_id","variable"])["value"]
    assert v[("andhra_pradesh","phonepe_count")] == 10862829
    assert v[("jammu_kashmir","phonepe_count")] == 110
    assert v[("jammu_kashmir","phonepe_amount_inr")] == 1050.0
    assert v[("jammu_kashmir","phonepe_users")] == 45
    assert "ladakh" not in set(df.state_id)
def test_missing_quarter_fails(tmp_path):
    with pytest.raises(FileNotFoundError, match="2018/2.json"):
        build_interim(FIX, tmp_path/"p.csv", ["2018Q1","2018Q2"])
\end{verbatim}
\item Run \f{.venv/bin/pytest tests/py/test\_phonepe.py -v}. Expected: FAIL (ImportError).
\item Implement. \f{build\_interim} maps names, then sums with \f{groupby(["state\_id","period","variable"])}.
      \f{fetch} runs: \f{git clone --filter=blob:none --no-checkout <url> <dest>/repo}; \f{git sparse-checkout set
      data/map/transaction/hover/country/india data/map/user/hover/country/india}; \f{git checkout <commit>};
      then appends a manifest row.
\item Run the tests. Expected: PASS.
\item Run on real data: \f{PYTHONPATH=src/py .venv/bin/python -m btpdata.fetch\_phonepe \&\& PYTHONPATH=src/py .venv/bin/python -m btpdata.parse\_phonepe}.
      Expected: \f{data/interim/phonepe.csv} has $35 \times 34 \times 3 = 3570$ rows.
\item Commit (code, tests, fixtures, \f{data/raw/phonepe/manifest.csv}, \f{data/interim/phonepe.csv}):
      \f{git commit -m "feat: PhonePe Pulse ingestion"}.
\end{steps}

% =====================================================================
\section{Task 3: Population}
\textbf{Files.} Create \f{data/manual/population\_ncp2020.csv}, \f{data/manual/manifest.csv},
\f{src/py/btpdata/population.py}, \f{tests/py/test\_population.py}.

\textbf{Interfaces.} Produces \f{load\_annual(path: Path) -> pd.DataFrame} (\f{state\_id, year, population},
in persons, merged ids summed); \f{to\_quarterly(annual: pd.DataFrame, periods: list[str]) -> pd.DataFrame}
(\f{state\_id, period, population}); \f{main()} writes \f{data/interim/population.csv} (variable \f{population}).

\textbf{Interpolation rule.} The value for year $y$ is the $y$Q1 value. For $k = 1,2,3$:
$p_{y\mathrm{Q}(k+1)} = p_y + \tfrac{k}{4}(p_{y+1} - p_y)$.

\begin{steps}
\item Transcribe the state table of projected total population (as of 1 March), years 2018--2027, from the
      \emph{Report of the Technical Group on Population Projections: India and States 2011--2036} (National
      Commission on Population, MoHFW, July 2020) into \f{population\_ncp2020.csv}, with columns
      \f{raw\_name,year,population\_thousands}. Include the India row as \f{raw\_name = "india"}. Record the table
      number and page in \f{manifest.csv}.
\item Write the failing tests:
\begin{verbatim}
def test_interp():
    a = pd.DataFrame({"state_id":["x"]*2,"year":[2018,2019],"population":[1000.,1400.]})
    q = to_quarterly(a, ["2018Q1","2018Q2","2018Q3","2018Q4","2019Q1"])
    assert q.population.tolist() == [1000.,1100.,1200.,1300.,1400.]
def test_merge(tmp_path):  # csv with jammu & kashmir 900, ladakh 100 (thousands), 2018
    assert load_annual(p).query("state_id=='jammu_kashmir'").population.item() == 1_000_000
def test_needs_next_year():
    with pytest.raises(ValueError, match="2027"): to_quarterly(annual_2018_to_2026, ["2026Q2"])
def test_real_file_sums():  # each year: |sum(states) / india - 1| < 0.01
def test_real_file_complete():  # to_quarterly(real, quarters()) has no NaN for the 35 ids
\end{verbatim}
\item Run \f{.venv/bin/pytest tests/py/test\_population.py -v}. Expected: FAIL.
\item Implement \f{population.py}. The India row is used only in the sum test; \f{load\_annual} drops it.
      If a small UT is not in the report, \f{test\_real\_file\_complete} fails. Then stop and ask the user
      (options: a second source, or leave the UT out of ALL\_IDS).
\item Run the tests. Expected: PASS.
\item Commit: \f{git commit -m "feat: population projections, quarterly"}.
\end{steps}

% =====================================================================
\section{Task 4: State CPI deflator}
\textbf{Files.} Create \f{data/raw/cpi/} (manual download + manifest), \f{src/py/btpdata/cpi.py},
\f{tests/py/test\_cpi.py}, \f{tests/py/fixtures/cpi/} (header + 3 rows copied from the real file, same format).

\textbf{Interfaces.} Produces \f{read\_raw(path: Path) -> pd.DataFrame} (\f{state\_id, month, cpi, base});
\f{link\_bases(df, factors: dict[str, float]) -> pd.DataFrame} (all values on base 2012=100);
\f{to\_quarterly(df) -> pd.DataFrame} (\f{state\_id, period, cpi}); \f{main()} writes \f{data/interim/cpi.csv},
with variables \f{cpi} and \f{cpi\_fallback} (1 if the all-India CPI was used).

\begin{steps}
\item Download from MoSPI (eSankhyiki portal, CPI): monthly state-wise CPI \emph{Combined, General index},
      January 2018 -- June 2026, in every base the portal gives. If there is a newer base, also get the official
      linking factor. Save the files to \f{data/raw/cpi/}, add manifest rows, and copy a fixture.
\item Write the failing tests:
\begin{verbatim}
# q(cpi=[...]): test-local helper; builds Jan..Mar 2018 rows for one state, returns to_quarterly(...).cpi[0]
def test_quarter_mean(): assert q(cpi=[100,101,102]) == 101.0      # Jan-Mar -> 2018Q1
def test_quarter_needs_three_months(): assert math.isnan(q(cpi=[100,101]))
def test_link_no_jump():  # old base Dec-2025 = 180.0; new base Jan-2026 = 92.5; factor 1.95
    s = link_bases(df, {"all_india":1.95}); assert abs(s.cpi.iloc[-1]/180.0 - 1) < 0.02
def test_fixture_reads(): df = read_raw(FIX); assert set(df.columns)=={"state_id","month","cpi","base"}
def test_fallback(): # a state absent from the file gets all-India values and cpi_fallback == 1
\end{verbatim}
\item Run \f{.venv/bin/pytest tests/py/test\_cpi.py -v}. Expected: FAIL.
\item Implement \f{cpi.py}. \f{read\_raw} is written to match the fixture's real format. For merged ids, use the
      CPI of the larger-population component (J\&K for \f{jammu\_kashmir}). If the file gives a per-state linking
      factor, use it; if not, use the all-India factor, and record which one in the manifest \f{notes}.
\item Run the tests. Expected: PASS.
\item Commit: \f{git commit -m "feat: state CPI deflator"}.
\end{steps}

% =====================================================================
\section{Task 5: Quarterly panel build and checks}
\textbf{Files.} Create \f{src/py/btpdata/build\_panels.py}, \f{tests/py/test\_build\_panels.py}.

\textbf{Interfaces.} Consumes the interim files of Tasks 2--4. Produces
\f{build\_quarterly(interim\_dir: Path) -> pd.DataFrame};
\f{check\_panel(df: pd.DataFrame) -> None} (raises \f{PanelCheckError});
\f{coverage(df: pd.DataFrame) -> pd.DataFrame} (\f{variable, state\_id, n\_missing});
\f{main()} writes \f{data/final/panel\_quarterly.csv} and \f{data/final/coverage\_report.csv}.

\textbf{Output columns.} \f{state\_id, period, in\_main\_sample, phonepe\_count, phonepe\_amount\_inr, phonepe\_users,
population, cpi, cpi\_fallback, count\_pc, real\_amount\_pc, users\_per\_1000, ln\_count\_pc, ln\_real\_amount\_pc}.
Here \f{count\_pc} = count/pop; \f{real\_amount\_pc} = amount$\times$100/cpi/pop; \f{users\_per\_1000} = users$\times$1000/pop.

\begin{steps}
\item Write the failing tests (build small interim CSVs in \f{tmp\_path}):
\begin{verbatim}
def test_merged_per_capita(): # jammu_kashmir count 110, population 1000 -> count_pc == 0.11
def test_duplicate():  with pytest.raises(PanelCheckError, match="duplicate"): check_panel(dup)
def test_negative():   with pytest.raises(PanelCheckError, match="negative"):  check_panel(neg)
def test_outlier():    # one state count_pc = 11 x median of others in 2018Q1
    with pytest.raises(PanelCheckError, match=r"outlier.*goa.*2018Q1"): check_panel(out)
def test_zero_is_nan_not_inf():
    df = build_quarterly(d_zero)
    assert df.ln_count_pc.isna().sum() == 1
    assert np.isfinite(df.ln_count_pc.dropna()).all()
    assert coverage(df).query("variable=='ln_count_pc'").n_missing.sum() == 1
def test_main_sample_flag(): assert df.in_main_sample.sum() == 30 * n_periods
\end{verbatim}
\item Run \f{.venv/bin/pytest tests/py/test\_build\_panels.py -v}. Expected: FAIL.
\item Implement \f{build\_panels.py}: pivot the interim tables to wide, merge on (\f{state\_id}, \f{period}),
      derive the columns, run \f{check\_panel}, write the files. The outlier check covers \f{count\_pc},
      \f{real\_amount\_pc} and \f{users\_per\_1000}.
\item Run the tests. Expected: PASS.
\item Run on real data: \f{PYTHONPATH=src/py .venv/bin/python -m btpdata.build\_panels}. Expected: 1190 rows ($35\times34$). If the
      outlier check fails, stop and report the state, period and value to the user.
\item Commit (code, tests, \f{data/final/*.csv}): \f{git commit -m "feat: quarterly panel build with checks"}.
\end{steps}

% =====================================================================
\section{Task 6: NPCI state data and proxy validation}
\textbf{Files.} Create \f{data/raw/npci/} (manual download + manifest), \f{src/py/btpdata/npci.py},
\f{src/py/btpdata/validation.py}, \f{tests/py/test\_npci.py}, \f{tests/py/test\_validation.py}, fixture copied from one real file.

\textbf{Interfaces.} Produces \f{UNATTRIBUTED\_LABELS: frozenset[str]} (set from the real file);
\f{read\_month(path: Path, month: str) -> pd.DataFrame} (\f{state\_id, month, npci\_volume, npci\_value\_inr};
attributed rows only); \f{to\_quarterly(df) -> pd.DataFrame} (sum of three months, else row dropped);
\f{compare(panel: pd.DataFrame, npci\_q: pd.DataFrame) -> pd.DataFrame} with columns
\f{period, measure, pearson, spearman, n\_states}, measures \f{share\_count, share\_value, count\_pc, value\_pc};
\f{main()} writes \f{data/final/validation\_npci.csv} for 2025Q2--2026Q2.

\begin{steps}
\item Download the NPCI state-wise UPI files, April 2025 -- June 2026, from the NPCI UPI ecosystem statistics page
      in a browser (the page needs JavaScript). Save them to \f{data/raw/npci/} with manifest rows, and record the
      units (for example Mn, Cr) in \f{notes}. Copy one file's header and 3 rows plus the unattributed row as the fixture.
\item Write the failing tests:
\begin{verbatim}
def test_unattributed_dropped():
    df = read_month(FIX, "2025-04"); assert not df.state_id.isin(["others"]).any()
def test_units(): # value converted to rupees using the unit in the header (e.g. Cr -> x 1e7)
def test_unknown_row_fails(): # a row not in crosswalk and not in UNATTRIBUTED_LABELS -> UnknownStateError
def test_shares_over_attributed():  # volumes 10,30,60 + Others 100 -> shares 0.1,0.3,0.6
def test_perfect_agreement():
    r = compare(p, n); assert (r.pearson.round(9) == 1).all() and (r.spearman == 1).all()
def test_n_states_intersection(): # 3 states in panel, 2 in npci -> n_states == 2
\end{verbatim}
\item Run \f{.venv/bin/pytest tests/py/test\_npci.py tests/py/test\_validation.py -v}. Expected: FAIL.
\item Implement \f{npci.py} (written to match the fixture's real format) and \f{validation.py} (pandas
      \f{corr(method="pearson")} and \f{corr(method="spearman")}, over the states present in both).
\item Run the tests. Expected: PASS.
\item Run on real data. Expected: 20 rows (5 quarters $\times$ 4 measures). If any \f{share\_count} Pearson
      value is below 0.8, report it to the user (spec risk~2). Do not change the method.
\item Commit: \f{git commit -m "feat: NPCI validation of PhonePe proxy"}.
\end{steps}

% =====================================================================
\section{Task 7: State boundaries for maps}
\textbf{Files.} Create \f{src/py/btpdata/boundaries.py}, \f{tests/py/test\_boundaries.py}, \f{data/raw/boundaries/}.

\textbf{Interfaces.} Produces \f{add\_state\_ids(geojson: dict, name\_field: str) -> dict} (adds property
\f{state\_id}; merged ids keep separate features with the same id); \f{main()} writes
\f{data/interim/india\_states.geojson}.

\begin{steps}
\item Download geoBoundaries \emph{gbOpen IND ADM1} (simplified GeoJSON, CC BY 4.0) at a fixed release into
      \f{data/raw/boundaries/}, and add a manifest row with the release tag.
\item Write the failing tests: two features named ``Orissa'' and ``Ladakh'' $\to$ \f{state\_id} \f{odisha} and
      \f{jammu\_kashmir}; a feature named ``Atlantis'' raises \f{UnknownStateError}; on the real file, the set of
      ids equals \f{set(ALL\_IDS)}.
\item Run \f{.venv/bin/pytest tests/py/test\_boundaries.py -v}. Expected: FAIL.
\item Implement \f{boundaries.py}. Add any new name variants to \f{state\_crosswalk.csv}.
\item Run the tests. Expected: PASS.
\item Commit: \f{git commit -m "feat: state boundaries with state ids"}.
\end{steps}

% =====================================================================
\section{Task 8: R environment and descriptive results}
\textbf{Files.} Create \f{src/R/lib/descriptives\_fns.R}, \f{src/R/01\_descriptives.R}, \f{tests/R/test\_descriptives.R},
\f{renv.lock}, \f{.Rprofile}.

\textbf{Interfaces.} Consumes \f{panel\_quarterly.csv}, \f{validation\_npci.csv}, \f{india\_states.geojson}. Produces
\f{sigma\_cv(dt, var) -> data.table(period, cv)} (cross-state \f{sd(n-1)/mean} for each period, main sample);
\f{beta\_convergence(dt, var, start = "2018Q4", end = "2026Q2") -> lm} (annualised growth $= (\ln x_{end} - \ln
x_{start})/7.5$ regressed on $\ln x_{start}$). The start is 2018Q4 to avoid the very small 2018Q1 base;
7.5 years separate 2018Q4 and 2026Q2.

\textbf{Outputs.} \f{output/figures/}: \f{adoption\_trend\_main.pdf} (median and IQR of \f{ln\_real\_amount\_pc}, main
sample), \f{adoption\_map\_2018Q4.pdf}, \f{adoption\_map\_2026Q2.pdf} (choropleth of \f{real\_amount\_pc}),
\f{sigma\_convergence.pdf}. \f{output/tables/}: \f{beta\_convergence.tex} (modelsummary, \f{output = "latex"}),
\f{validation\_npci.tex}, \f{summary\_stats.tex}.

\begin{steps}
\item Initialise renv: \f{Rscript -e 'renv::init(bare=TRUE)'}; then
      \f{Rscript -e 'install.packages(c("data.table","ggplot2","sf","modelsummary","testthat"))'};
      then \f{Rscript -e 'renv::snapshot()'}.
\item Write the failing tests:
\begin{verbatim}
test_that("cv", expect_equal(sigma_cv(dt3, "x")$cv, 0.5))          # values 1,2,3 in one period
test_that("beta", expect_equal(unname(coef(beta_convergence(dts, "x"))[2]), -0.02, tolerance=1e-9))
# dts built so that growth == 0.1 - 0.02 * ln(x_start) exactly
\end{verbatim}
\item Run \f{Rscript -e 'testthat::test\_dir("tests/R")'}. Expected: FAIL.
\item Implement \f{descriptives\_fns.R}, then \f{01\_descriptives.R}, which sources it and writes all the outputs.
\item Run the tests (expected: PASS), then \f{Rscript src/R/01\_descriptives.R}. Expected: all 7 output files exist
      and are not empty.
\item Commit (code, tests, \f{renv.lock}, \f{output/}): \f{git commit -m "feat: descriptive results"}.
\end{steps}

% =====================================================================
\section{Task 9: Makefile and end-to-end run}
\textbf{Files.} Create \f{Makefile}.

\textbf{Targets.} \f{test} (pytest + testthat); \f{data} (Tasks 2--7 \f{main()}s, in order); \f{analysis}
(\f{01\_descriptives.R}); \f{report} (\f{latexmk} in \f{thesis/btp1} and \f{thesis/btp1-slides});
\f{all} = \f{test data analysis report}. Manual downloads are prerequisites: \f{data} stops with a clear
message if \f{data/raw/cpi/} or \f{data/raw/npci/} is empty.

\begin{steps}
\item Write the Makefile.
\item Run \f{make test data analysis} from a clean clone in the session scratchpad
      (\f{git clone . \$SCRATCH/e2e}, where \f{\$SCRATCH} is the scratchpad directory of the executing session),
      and copy \f{data/raw/} into it. Expected: exit 0, and the \f{data/final/} and \f{output/} files match the committed
      ones (\f{git status} clean).
\item Read \f{data/final/coverage\_report.csv} and list every main-sample variable with missing cells for the user.
\item Commit: \f{git commit -m "build: Makefile pipeline"}.
\end{steps}

% =====================================================================
\section{Task 10: BTP-I report (\LaTeX)}
\textbf{Files.} Create \f{thesis/btp1/main.tex}, \f{thesis/btp1/chapters/ch1\_intro.tex} \ldots
\f{ch6\_plan.tex}, \f{thesis/btp1/refs.bib}.

\textbf{Setup.} \f{report} class, 12pt, A4 (or the department template from Task~0); packages \f{natbib}
(\f{plainnat}), \f{graphicx}, \f{booktabs}, \f{hyperref}; \f{\textbackslash graphicspath\{\{../../output/figures/\}\}};
tables via \f{\textbackslash input\{../../output/tables/<name>.tex\}}.

\textbf{refs.bib.} The seven key references in the spec, plus: ICRIER WP IPCIDE-wp1 (Reddy, Kedia, Shukla); Singh \&
Ghosh (2021), \emph{Economic Analysis and Policy} 71:674--693; the PhonePe Pulse dataset (\f{@misc}, with the commit hash);
NPCI UPI statistics (\f{@misc}, with the access date); the NCP 2020 population report.

\textbf{Chapter outlines} (as \f{\%} comments at the top of each chapter file, from spec Section~6):
Ch~1 background, problem, RQ1--RQ4 and hypotheses H1--H4; Ch~2 finance--growth, FI--ICT (Pradhan et al.),
digital payments, gap; Ch~3 data availability findings, PhonePe as proxy, variable table, sample, the
\f{coverage\_report} summary; Ch~4 Steps 0--3 of the spec; Ch~5 \f{adoption\_trend\_main}, the maps,
\f{sigma\_convergence}, \f{beta\_convergence}, \f{validation\_npci}, \f{summary\_stats}; Ch~6 the BTP-II plan.

\begin{steps}
\item Create the files. Each chapter has its outline comments, the section headings, and the figure and table
      includes listed above.
\item Run \f{cd thesis/btp1 \&\& latexmk -pdf main.tex}. Expected: exit 0; \f{grep -E "undefined|Overfull" main.log}
      prints nothing.
\item Commit: \f{git commit -m "docs: BTP-I report skeleton"}.
\item The student writes the chapters (weeks 3--4). Run the same compile check after each chapter.
\end{steps}

% =====================================================================
\section{Task 11: BTP-I slides (Beamer)}
\textbf{Files.} Create \f{thesis/btp1-slides/slides.tex} (shares \f{../btp1/refs.bib}).

\textbf{Frames (14).} Title; Motivation; RQs and hypotheses; Literature and gap; Data availability; PhonePe Pulse
as proxy; Variables and sample; Method plan; Adoption trend; Maps 2018Q4 vs 2026Q2; Convergence; NPCI
validation; BTP-II plan; References.

\begin{steps}
\item Create \f{slides.tex} with the 14 frames, the figure includes and the speaker notes as \f{\%} comments.
\item Run \f{cd thesis/btp1-slides \&\& latexmk -pdf slides.tex}. Expected: exit 0;
      \f{grep -c '\textbackslash begin\{frame\}' slides.tex} prints 14.
\item Commit: \f{git commit -m "docs: BTP-I slides"}.
\end{steps}

\end{document}
